% ============================================================================
% Chapter 9 solutions: Simplex Method in Tableau Form
% Book source: part1-linear-programming/ch06-simplex/simplex-tableau.tex
% Exercise numbering follows \begin{ex} order in that file (12 exercises).
% All pivots and optima verified with exact-fraction basis algebra and
% scipy.optimize.linprog, July 2026.
% ============================================================================
\section{Bab 9: Metode Simpleks dalam Bentuk Tableau}

Bab 9 menyajikan kembali metode simpleks sebagai operasi baris pada satu
tableau. Kedua belas latihannya mencakup pembacaan tableau (9.1 sampai 9.3),
pelaksanaan pivot dan penyelesaian lengkap, termasuk penyusunan Big-M (9.4 sampai
9.7), konvensi tanda dan korespondensi tableau--kamus (9.8, 9.9), pengenalan
ketakterbatasan dan ketaklayakan dari sebuah tableau (9.10, 9.11), serta satu
soal tantangan yang mengubah tableau tak terbatas menjadi sinar sertifikat yang
eksplisit (9.12). Ingat selalu konvensi tanda bab ini: baris $z$ menyimpan
$z - \c^\top\x = 0$, sehingga variabel masuk memiliki entri baris $z$ yang
\emph{negatif}, dan tableau optimal ketika setiap entri baris $z$ tak negatif.
Semua aritmetika di bawah ini diverifikasi dengan aljabar basis menggunakan
pecahan eksak, dan setiap optimum dikonfirmasi dengan pemecah program linier. Solusi
pilihan dalam buku tersedia untuk 9.2, 9.3, 9.5, dan 9.6 serta direproduksi di
sini.

\exsol{9.1}{Membaca tableau di tengah proses penyelesaian}{1}

\textbf{(a)} Variabel basis adalah label baris: $s_1, s_2, y$. Variabel
nonbasis adalah $x$ dan $s_3$.

\textbf{(b)} Tetapkan $x = s_3 = 0$, lalu baca kolom RHS:
\[
x = 0,\quad y = 7,\quad s_1 = 2,\quad s_2 = 9,\quad s_3 = 0,
\qquad z = 21.
\]

\textbf{(c)} Entri baris $z$ untuk variabel nonbasis adalah
$-\tfrac12$ (untuk $x$) dan $\tfrac32$ (untuk $s_3$). Dalam konvensi bab
ini, keduanya adalah biaya tereduksi yang dicatat; karena tandanya dibalik saat
disimpan, menaikkan $x$ satu unit akan \emph{menaikkan} $z$ sebesar $\tfrac12$,
sedangkan menaikkan $s_3$ akan menurunkannya sebesar $\tfrac32$.

\textbf{(d)} Belum optimal: baris $z$ memiliki entri negatif $-\tfrac12$
di bawah $x$, sehingga $x$ masuk berikutnya. (Uji rasio sesudahnya memberikan
$2/\tfrac12 = 4$, $9/\tfrac32 = 6$, $7/\tfrac12 = 14$, sehingga $s_1$ akan
keluar dan menghasilkan tableau optimal dengan $z = 23$.)

\exsol{9.2}{Membaca tableau}{1}

(Solusi pilihan dalam buku, telah diverifikasi.) Variabel basis adalah label
baris $x_1$ dan $x_2$; variabel nonbasis adalah $s_1$ dan $s_2$. Menetapkan
variabel nonbasis sama dengan nol dan membaca kolom RHS memberikan solusi basis
layak
\[
x_1 = 6, \quad x_2 = 4, \quad s_1 = s_2 = 0, \qquad z = 40.
\]
Entri baris objektif untuk variabel nonbasis adalah $3$ dan $1$, keduanya tak
negatif. Dalam konvensi ini, entri positif berarti bahwa menaikkan variabel
nonbasis yang bersangkutan akan \emph{menurunkan} $z$, sehingga tidak ada pivot
yang dapat memperbaiki objektif dan tableau tersebut optimal.

\exsol{9.3}{Variabel masuk, uji rasio, dan elemen pivot}{1}

(Solusi pilihan dalam buku, telah diverifikasi.) Entri baris objektif yang
paling negatif adalah $-4$, sehingga $x_1$ masuk. Uji rasio membagi setiap RHS
dengan entri positif pada kolom $x_1$:
\[
\frac{6}{2} = 3, \qquad \frac{8}{4} = 2, \qquad \frac{12}{3} = 4.
\]
Rasio minimum $2$ terjadi pada baris $s_2$, sehingga $s_2$ keluar dan elemen
pivotnya adalah $4$ pada baris $s_2$ di kolom $x_1$.

\exsol{9.4}{Melakukan satu pivot}{2}

Lakukan pivot pada $4$ di baris $s_2$, kolom $x_1$. Kalikan baris pivot dengan
$\tfrac14$ untuk memperoleh baris $x_1$ yang baru,
$[\,1,\ \tfrac14,\ 0,\ \tfrac14,\ 0 \mid 2\,]$, lalu nolkan kolom $x_1$:
$\text{baris }s_1 \leftarrow \text{baris }s_1 - 2(\text{baris }x_1)$,
$\text{baris }s_3 \leftarrow \text{baris }s_3 - 3(\text{baris }x_1)$,
$\text{baris }z \leftarrow \text{baris }z + 4(\text{baris }x_1)$. Hasilnya adalah:
\begin{center}
\renewcommand{\arraystretch}{1.3}
\begin{tabular}{c|ccccc|c}
Basis & $x_1$ & $x_2$ & $s_1$ & $s_2$ & $s_3$ & RHS\\
\hline
$z$   & $0$ & $-2$ & $0$ & $1$ & $0$ & $8$\\
\hline
$s_1$ & $0$ & $\tfrac52$  & $1$ & $-\tfrac12$ & $0$ & $2$\\
$x_1$ & $1$ & $\tfrac14$  & $0$ & $\tfrac14$  & $0$ & $2$\\
$s_3$ & $0$ & $\tfrac{13}{4}$ & $0$ & $-\tfrac34$ & $1$ & $6$\\
\end{tabular}
\end{center}
Kolom $x_1$ sekarang merupakan kolom satuan $(0; 0, 1, 0)$, dan nilai
objektif pada RHS baris $z$ telah meningkat dari $0$ menjadi $8$ (solusi basis
layak yang baru adalah $x_1 = 2$, $x_2 = 0$ dengan
$z = 4 \cdot 2 = 8$). Tableau tersebut belum optimal ($-2$ di bawah $x_2$),
tetapi latihan hanya meminta satu pivot ini.

\exsol{9.5}{Menyelesaikan program linier dalam bentuk tableau}{2}

(Solusi pilihan dalam buku, telah diverifikasi, dengan tableau awal
ditampilkan.) Setelah menambahkan variabel slack $s_1, s_2$, tableau awalnya
adalah
\begin{center}
\renewcommand{\arraystretch}{1.3}
\begin{tabular}{c|cccc|c|c}
Basis & $x$ & $y$ & $s_1$ & $s_2$ & RHS & Rasio\\
\hline
$z$   & $-1$ & $\mathbf{-2}$ & $0$ & $0$ & $0$ & \\
\hline
$s_1$ & $1$ & $1$ & $1$ & $0$ & $4$ & $4/1 = 4$\\
$s_2$ & $1$ & $\fbox{3}$ & $0$ & $1$ & $6$ & $6/3 = 2$ (min)\\
\end{tabular}
\end{center}

\textbf{Pivot 1:} $y$ masuk ($-2$ adalah yang paling negatif); rasionya $4$
dan $2$, sehingga $s_2$ keluar; lakukan pivot pada $3$ yang diberi kotak.
Membagi baris pivot dengan $3$ dan menolkan kolom $y$ memberikan
\begin{center}
\renewcommand{\arraystretch}{1.3}
\begin{tabular}{c|cccc|c|c}
Basis & $x$ & $y$ & $s_1$ & $s_2$ & RHS & Rasio\\
\hline
$z$   & $\mathbf{-\tfrac13}$ & $0$ & $0$ & $\tfrac23$ & $4$ & \\
\hline
$s_1$ & $\fbox{$\tfrac23$}$ & $0$ & $1$ & $-\tfrac13$ & $2$ & $2/\tfrac23 = 3$ (min) \\
$y$   & $\tfrac13$ & $1$ & $0$ & $\tfrac13$  & $2$ & $2/\tfrac13 = 6$\\
\end{tabular}
\end{center}

\textbf{Pivot 2:} $x$ masuk ($-\tfrac13$); rasionya $3$ dan $6$, sehingga
$s_1$ keluar; lakukan pivot pada $\tfrac23$ yang diberi kotak:
\begin{center}
\renewcommand{\arraystretch}{1.3}
\begin{tabular}{c|cccc|c}
Basis & $x$ & $y$ & $s_1$ & $s_2$ & RHS\\
\hline
$z$   & $0$ & $0$ & $\tfrac12$ & $\tfrac12$ & $5$\\
\hline
$x$ & $1$ & $0$ & $\tfrac32$  & $-\tfrac12$ & $3$\\
$y$ & $0$ & $1$ & $-\tfrac12$ & $\tfrac12$  & $1$\\
\end{tabular}
\end{center}
Setiap entri baris $z$ tak negatif, sehingga tableau tersebut optimal:
$(x, y) = (3, 1)$ dengan $z^* = 5$ (kedua kendala aktif).

\exsol{9.6}{Penyelesaian lengkap dengan dua pivot}{2}

(Solusi pilihan dalam buku, telah diverifikasi.) Setelah menambahkan variabel
slack $s_1, s_2$, tableau awalnya adalah
\begin{center}
\renewcommand{\arraystretch}{1.3}
\begin{tabular}{c|cccc|c|c}
Basis & $x$ & $y$ & $s_1$ & $s_2$ & RHS & Rasio\\
\hline
$z$   & $\mathbf{-3}$ & $-2$ & $0$ & $0$ & $0$ & \\
\hline
$s_1$ & $\fbox{2}$ & $1$ & $1$ & $0$ & $10$ & $10/2 = 5$ (min)\\
$s_2$ & $1$ & $3$ & $0$ & $1$ & $15$ & $15/1 = 15$\\
\end{tabular}
\end{center}

\textbf{Pivot 1:} $x$ masuk ($-3$ adalah yang paling negatif); $s_1$ keluar;
lakukan pivot pada $2$ yang diberi kotak. Kalikan baris pivot dengan
$\tfrac12$, lalu nolkan kolom $x$:
\begin{center}
\renewcommand{\arraystretch}{1.3}
\begin{tabular}{c|cccc|c|c}
Basis & $x$ & $y$ & $s_1$ & $s_2$ & RHS & Rasio\\
\hline
$z$   & $0$ & $\mathbf{-\tfrac12}$ & $\tfrac32$ & $0$ & $15$ & \\
\hline
$x$   & $1$ & $\tfrac12$ & $\tfrac12$ & $0$ & $5$ & $5/\tfrac12 = 10$\\
$s_2$ & $0$ & $\fbox{$\tfrac52$}$ & $-\tfrac12$ & $1$ & $10$ & $10/\tfrac52 = 4$ (min) \\
\end{tabular}
\end{center}

\textbf{Pivot 2:} $y$ masuk ($-\tfrac12$); rasionya $10$ dan $4$, sehingga
$s_2$ keluar; lakukan pivot pada $\tfrac52$ yang diberi kotak:
\begin{center}
\renewcommand{\arraystretch}{1.3}
\begin{tabular}{c|cccc|c}
Basis & $x$ & $y$ & $s_1$ & $s_2$ & RHS\\
\hline
$z$   & $0$ & $0$ & $\tfrac75$ & $\tfrac15$ & $17$\\
\hline
$x$ & $1$ & $0$ & $\tfrac35$ & $-\tfrac15$ & $3$\\
$y$ & $0$ & $1$ & $-\tfrac15$ & $\tfrac25$ & $4$\\
\end{tabular}
\end{center}
Setiap entri baris $z$ tak negatif, sehingga tableau tersebut optimal:
$(x, y) = (3, 4)$ dengan $z^* = 17$ dan $s_1 = s_2 = 0$ (kedua kendala
aktif pada optimum).

\exsol{9.7}{Menyusun tableau Big-M}{2}

Kendala yang diperluas adalah
\[
x + y - e_1 + a_1 = 4,
\qquad
x + 2y + s_2 = 10,
\qquad
x, y, e_1, s_2, a_1 \geq 0,
\]
dengan objektif $z - 4x - 3y + M a_1 = 0$. Jika dituliskan secara langsung,
tableau tersebut adalah
\begin{center}
\renewcommand{\arraystretch}{1.3}
\begin{tabular}{c|ccccc|c}
Basis & $x$ & $y$ & $e_1$ & $s_2$ & $a_1$ & RHS\\
\hline
$z$   & $-4$ & $-3$ & $0$ & $0$ & $M$ & $0$\\
\hline
$a_1$ & $1$ & $1$ & $-1$ & $0$ & $1$ & $4$\\
$s_2$ & $1$ & $2$ & $0$  & $1$ & $0$ & $10$\\
\end{tabular}
\end{center}
Ini belum merupakan tableau yang benar: variabel \emph{basis} $a_1$ memiliki
entri tak nol ($M$) pada baris $z$. Menolkan entri tersebut
($\text{baris }z \leftarrow \text{baris }z - M \cdot \text{baris }a_1$) memberikan
\begin{center}
\renewcommand{\arraystretch}{1.3}
\begin{tabular}{c|ccccc|c|c}
Basis & $x$ & $y$ & $e_1$ & $s_2$ & $a_1$ & RHS & Rasio\\
\hline
$z$   & $\mathbf{-4-M}$ & $-3-M$ & $M$ & $0$ & $0$ & $-4M$ & \\
\hline
$a_1$ & $\fbox{1}$ & $1$ & $-1$ & $0$ & $1$ & $4$ & $4/1 = 4$ (min) \\
$s_2$ & $1$ & $2$ & $0$  & $1$ & $0$ & $10$ & $10/1 = 10$\\
\end{tabular}
\end{center}
Sekarang entri objektif $a_1$ bernilai nol, sebagaimana diperlukan. Entri baris
$z$ yang paling negatif adalah $-4 - M$ (untuk $M$ besar, nilainya lebih kecil
daripada $-3 - M$), sehingga \emph{$x$ masuk terlebih dahulu}. Uji rasio
memberikan $4/1 = 4$ pada baris $a_1$, dibandingkan dengan $10/1 = 10$ pada
baris $s_2$, sehingga \emph{ya: variabel artifisial $a_1$ keluar pada pivot
pertama}, dan metode dilanjutkan dengan tableau layak biasa (pivot tersebut
menghasilkan $(x, y) = (4, 0)$ dengan $z = 16$).

\exsol{9.8}{Konvensi tanda pada baris $z$}{2}

\textbf{(a)} Sebuah kamus menampilkan
$z = \bar z + \sum_j \bar c_j x_j$ dalam variabel nonbasis. Tableau justru
menyimpan persamaan ini dengan semua suku di ruas kiri:
$z - \sum_j \bar c_j x_j = \bar z$. Setiap koefisien $\bar c_j$ berpindah
melewati tanda sama dengan sehingga tandanya berbalik. Karena itu, variabel yang
menarik dalam kamus ($\bar c_j$ positif) muncul dalam tableau dengan entri
negatif $-\bar c_j$. Dengan demikian, aturan masuk ``koefisien kamus paling
positif'' dan ``entri baris $z$ paling negatif'' memilih variabel yang sama.

\textbf{(b)} Dari $z = 12 + 3x_1 - 2x_2$, pindahkan suku-suku variabel ke kiri:
$z - 3x_1 + 2x_2 = 12$. Baris $z$ adalah
\[
\begin{array}{c|cc|c}
 & x_1 & x_2 & \text{RHS} \\
\hline
z & -3 & 2 & 12
\end{array}
\]
(Entri $0$ di bawah kolom slack atau kolom basis tidak ditampilkan.)

\textbf{(c)} Persamaan yang disimpan adalah
$z - \sum_j \bar c_j x_j = \bar z$, dan selain $z$, persamaan ini hanya
melibatkan variabel \emph{nonbasis}. Pada solusi basis saat ini, setiap variabel
nonbasis sama dengan nol, sehingga persamaan menyusut menjadi $z = \bar z$:
entri RHS adalah nilai objektif saat ini. Suku konstanta tidak pernah berpindah
melewati tanda sama dengan (suku itu sudah berada di kanan), sehingga inilah
satu-satunya bilangan pada baris tersebut yang tandanya tidak dibalik.

\exsol{9.9}{Dari tableau ke kamus}{2}

Setiap baris tableau merupakan persamaan. Baris $x_1$ menyatakan
$x_1 + 2s_1 - s_2 = 6$ dan baris $x_2$ menyatakan $x_2 - s_1 + s_2 = 4$;
baris $z$ menyatakan $z + 3s_1 + s_2 = 40$. Dengan menyelesaikan setiap
persamaan terhadap variabel basisnya (dan terhadap $z$), diperoleh:
\[
\begin{aligned}
\max\quad & z = 40 - 3s_1 - s_2 \\
\st\quad & x_1 = 6 - 2s_1 + s_2, \\
& x_2 = 4 + s_1 - s_2, \\
& x_1, x_2, s_1, s_2 \geq 0.
\end{aligned}
\]
Menetapkan variabel nonbasis $s_1 = s_2 = 0$ menghasilkan $x_1 = 6$,
$x_2 = 4$, $z = 40$: solusi basis layak yang sama seperti pada tableau.
(Kedua koefisien kamus dalam objektif bernilai negatif, sesuai dengan
optimalitas tableau.)

\exsol{9.10}{Mengenali masalah tak terbatas}{2}

Satu-satunya entri negatif pada baris $z$ adalah $-2$ di bawah $x_1$, sehingga
$x_1$ harus masuk. Uji rasio mencari entri \emph{positif} pada kolom $x_1$,
tetapi entrinya adalah $-1$ (pada baris $x_2$) dan $-3$ (pada baris $s_1$):
tidak ada entri positif, sehingga uji rasio gagal; tidak ada baris yang membatasi
kenaikan $x_1$. Menaikkan $x_1 = t$ sambil mempertahankan $s_2 = 0$ memberikan
$x_2 = 3 + t$ dan $s_1 = 5 + 3t$, yang tetap tidak negatif untuk semua
$t \geq 0$, sedangkan $z = 10 + 2t \to \infty$. Program linier tersebut
\textbf{tak terbatas}: program ini memiliki solusi layak dengan nilai objektif
sebesar apa pun dan tidak memiliki solusi optimal.

\exsol{9.11}{Membaca ketaklayakan}{2}

Ya, tableau tersebut optimal untuk masalah berpenalti (Big-M): untuk $M$ besar,
setiap entri baris $z$ ($0, 0, M, 2, 0$) tak negatif, sehingga tidak ada pivot
yang dapat memperbaiki objektif berpenalti. Namun, variabel artifisial masih
menjadi variabel basis dengan nilai positif $a_1 = 3$ (dibaca dari RHS). Solusi
layak bagi masalah asli harus memungkinkan semua variabel artifisial ditetapkan
sama dengan nol. Jadi, jika hasil terbaik masalah berpenalti masih mempertahankan
$a_1 = 3 > 0$, tidak ada solusi layak bagi program asli: program linier asli
tersebut \textbf{tak layak}.

Nilai objektif $8 - 3M$ menyampaikan hal yang sama: nilai ini memuat penalti
yang tidak dapat dihindari, $-M a_1 = -3M$. Karena $M$ menyatakan konstanta
yang besarnya sembarang, $8 - 3M$ dapat menjadi seburuk apa pun; masalah
berpenalti tidak dapat menghindari penalti tersebut. Inilah tepatnya ciri
aljabar dari ketaklayakan.

\exsol{9.12}{Menyertifikasi ketakterbatasan dari tableau}{3}

\textbf{(a)} Satu-satunya entri negatif pada baris $z$ adalah $-2$ di bawah
$x_2$, sehingga $x_2$ masuk. Entri kolomnya adalah $-1$ (pada baris $x_1$)
dan $0$ (pada baris $s_2$); uji rasio memerlukan entri \emph{positif}, tetapi
tidak menemukannya. Karena itu, tidak ada baris yang membatasi kenaikan $x_2$:
masalah tersebut tak terbatas.

\textbf{(b)} Tetapkan $x_2 = t \geq 0$ dan pertahankan variabel nonbasis
lainnya, $s_1 = 0$. Baris $x_1$ ($x_1 - x_2 + s_1 = 1$) memberikan
$x_1 = 1 + t$,
dan baris $s_2$ ($s_1 + s_2 = 3$) memberikan $s_2 = 3$. Jadi
\[
\x(t) = (x_1, x_2, s_1, s_2) = (1 + t,\ t,\ 0,\ 3), \qquad t \geq 0.
\]
Substitusi langsung ke dalam kendala asli memberikan:
\[
x_1 - x_2 = (1 + t) - t = 1 \leq 1 \quad\checkmark
\qquad
-x_1 + x_2 = -(1 + t) + t = -1 \leq 2 \quad\checkmark
\]
dan $x_1 = 1 + t \geq 0$, $x_2 = t \geq 0$, $s_1 = 0 \geq 0$,
$s_2 = 2 - (-1) = 3 \geq 0$. Jadi, $\x(t)$ layak untuk setiap
$t \geq 0$.

\textbf{(c)} Sepanjang sinar tersebut,
$z(t) = x_1 + x_2 = 1 + 2t \to \infty$ ketika $t \to \infty$. Secara
geometris, daerah layak adalah pita di antara dua garis sejajar
$x_1 - x_2 = 1$ dan $-x_1 + x_2 = 2$ pada kuadran pertama; daerah ini tak
terbatas dalam arah $\mathbf{d} = (1, 1)$. Sinar sertifikat berawal di titik
sudut $(1, 0)$ dan mengarah sepanjang $(1, 1)$, mengikuti batas
$x_1 - x_2 = 1$; karena $\c^\top \mathbf{d} = 1 + 1 = 2 > 0$, objektif tumbuh
tanpa batas di sepanjang sinar itu.
